% GENERATED FILE. Edit canonical lessons, then run npm run generate:books.
% canonical-source-hash: fnv1a64:f851b57180f1cb6a
% canonical-lessons: JA-C337-setsuyaku-suru, JA-C337-denryoku, JA-C337-nenryou, JA-C337-sekiyu, JA-C337-mottainai

\chapter{Saving Energy}
\label{ch:ja-saving-energy}
\hlchaptermodality{337}

\begin{chapteropening}
\textbf{By the end of this chapter:} \emph{I can say how we save energy and why we should.}

Complete the last lesson of chapter 337: I can say how we save energy and why we should.
\end{chapteropening}

\section[\texorpdfstring{\ja{せつやくする} (setsuyaku suru)}{せつやくする (setsuyaku suru)}]{\ja{せつやくする} (setsuyaku suru) — to economize; to save (energy, money)}

\label{lesson:JA-C337-setsuyaku-suru}



\begin{quote}
Before the new one: say the Japanese for air, then the Japanese for pollution.
\end{quote}

\subsection*{You'll want to know: \ja{せつやくする}}
\textbf{\ja{せつやくする}} — \emph{setsuyaku suru} — \textquotedblleft{}to economize; to save (energy, money)\textquotedblright{}.

\begin{quote}
\textbf{\ja{でんきが} \ja{たかいので、せつやくしています。}} — \emph{denki ga takai node, setsuyaku shite imasu} — \textquotedblleft{}Electricity is expensive, so we are cutting down.\textquotedblright{}
\end{quote}

Polite \textbf{\ja{せつやくします}}, \emph{setsuyakushimasu}; a \ja{する} verb. Past (polite): \textbf{\ja{ははが} \ja{みずを} \ja{せつやくしました。}}, \emph{haha ga mizu o setsuyaku shimashita}.

\begin{cousinweb}
\textbf{\ja{せつやくする}} (\emph{setsuyaku suru}) is built from a Sino-Japanese noun and the native verb \emph{suru}, \textquotedblleft{}to do\textquotedblright{}. The noun, \emph{setsuyaku}, is two Chinese characters read with their on'yomi, readings borrowed from Chinese: \emph{setsu}, \textquotedblleft{}joint\textquotedblright{} (Middle Chinese \emph{tset}) and \emph{yaku}, \textquotedblleft{}shrink\textquotedblright{} (\emph{'jak}) (KANJIDIC2; JMdict; Middle Chinese in Baxter's transcription).
\end{cousinweb}

\begin{grammarlens}[title={What you've built}]
One more word for giving reasons and opinions.
\end{grammarlens}

\subsection*{Your turn}
\begin{itemize}
\raggedright
  \item \emph{Say it:} \emph{setsuyaku suru}
  \item \emph{Say it:} \emph{setsuyaku suru}, once more
  \item \emph{Say it:} \emph{denki ga takai node, setsuyaku shite imasu} — Electricity is expensive, so we are cutting down.
  \item \emph{Recall:} say the Japanese for air, then the Japanese for pollution, then say \emph{setsuyaku suru} again
\end{itemize}

\begin{tcolorbox}[breakable,colback=teal!4,colframe=teal!35!black,title={Before you move on}]
What does \ja{せつやくする} mean? (\textquotedblleft{}To economize; to save (energy, money)\textquotedblright{}.) Say it once more.
\end{tcolorbox}
\section[\texorpdfstring{\ja{でんりょく} (denryoku)}{でんりょく (denryoku)}]{\ja{でんりょく} (denryoku) — electric power}

\label{lesson:JA-C337-denryoku}



\begin{quote}
Before the new one: say the Japanese for pollution, then the Japanese for to economize.
\end{quote}

\subsection*{You'll want to know: \ja{でんりょく}}
\textbf{\ja{でんりょく}} — \emph{denryoku} — \textquotedblleft{}electric power\textquotedblright{}.

\begin{quote}
\textbf{\ja{なつは} \ja{でんりょくが} \ja{たりないと} \ja{おもいます。}} — \emph{natsu wa denryoku ga tarinai to omoimasu} — \textquotedblleft{}I think there is not enough electric power in summer.\textquotedblright{}
\end{quote}

\begin{cousinweb}
\textbf{\ja{でんりょく}} (\emph{denryoku}) is Sino-Japanese (kango): two Chinese characters read with their on'yomi, readings borrowed from Chinese. They are \emph{den}, \textquotedblleft{}electricity\textquotedblright{} (Middle Chinese \emph{denH}) and \emph{ryoku}, \textquotedblleft{}power\textquotedblright{} (\emph{lik}) (KANJIDIC2; Middle Chinese in Baxter's transcription).
\end{cousinweb}

\begin{grammarlens}[title={What you've built}]
One more word for giving reasons and opinions.
\end{grammarlens}

\subsection*{Your turn}
\begin{itemize}
\raggedright
  \item \emph{Say it:} \emph{denryoku}
  \item \emph{Say it:} \emph{denryoku}, once more
  \item \emph{Say it:} \emph{natsu wa denryoku ga tarinai to omoimasu} — I think there is not enough electric power in summer.
  \item \emph{Recall:} say the Japanese for pollution, then the Japanese for to economize, then say \emph{denryoku} again
\end{itemize}

\begin{tcolorbox}[breakable,colback=teal!4,colframe=teal!35!black,title={Before you move on}]
What does \ja{でんりょく} mean? (\textquotedblleft{}Electric power\textquotedblright{}.) Say it once more.
\end{tcolorbox}
\section[\texorpdfstring{\ja{ねんりょう} (nenryō)}{ねんりょう (nenryō)}]{\ja{ねんりょう} (nenryō) — fuel}

\label{lesson:JA-C337-nenryou}



\begin{quote}
Before the new one: say the Japanese for to economize, then the Japanese for electric power.
\end{quote}

\subsection*{You'll want to know: \ja{ねんりょう}}
\textbf{\ja{ねんりょう}} — \emph{nenryō} — \textquotedblleft{}fuel\textquotedblright{}.

\begin{quote}
\textbf{\ja{ねんりょうが} \ja{たかく} \ja{なったので、こまっています。}} — \emph{nenryō ga takaku natta node, komatte imasu} — \textquotedblleft{}Fuel has become expensive, so we are in trouble.\textquotedblright{}
\end{quote}

\begin{cousinweb}
\textbf{\ja{ねんりょう}} (\emph{nenryō}) is Sino-Japanese (kango): two Chinese characters read with their on'yomi, readings borrowed from Chinese. They are \emph{nen}, \textquotedblleft{}burn\textquotedblright{} and \emph{ryō}, \textquotedblleft{}materials\textquotedblright{} (Middle Chinese \emph{lew}) (KANJIDIC2; Middle Chinese in Baxter's transcription).
\end{cousinweb}

\begin{grammarlens}[title={What you've built}]
One more word for giving reasons and opinions.
\end{grammarlens}

\subsection*{Your turn}
\begin{itemize}
\raggedright
  \item \emph{Say it:} \emph{nenryō}
  \item \emph{Say it:} \emph{nenryō}, once more
  \item \emph{Say it:} \emph{nenryō ga takaku natta node, komatte imasu} — Fuel has become expensive, so we are in trouble.
  \item \emph{Recall:} say the Japanese for to economize, then the Japanese for electric power, then say \emph{nenryō} again
\end{itemize}

\begin{tcolorbox}[breakable,colback=teal!4,colframe=teal!35!black,title={Before you move on}]
What does \ja{ねんりょう} mean? (\textquotedblleft{}Fuel\textquotedblright{}.) Say it once more.
\end{tcolorbox}
\section[\texorpdfstring{\ja{せきゆ} (sekiyu)}{せきゆ (sekiyu)}]{\ja{せきゆ} (sekiyu) — petroleum; oil}

\label{lesson:JA-C337-sekiyu}



\begin{quote}
Before the new one: say the Japanese for electric power, then the Japanese for fuel.
\end{quote}

\subsection*{You'll want to know: \ja{せきゆ}}
\textbf{\ja{せきゆ}} — \emph{sekiyu} — \textquotedblleft{}petroleum; oil\textquotedblright{}.

\begin{quote}
\textbf{\ja{せきゆは} \ja{すくなく} \ja{なると} \ja{おもいます。}} — \emph{sekiyu wa sukunaku naru to omoimasu} — \textquotedblleft{}I think oil will become scarce.\textquotedblright{}
\end{quote}

\begin{cousinweb}
\textbf{\ja{せきゆ}} (\emph{sekiyu}) is Sino-Japanese (kango): two Chinese characters read with their on'yomi, readings borrowed from Chinese. They are \emph{seki}, \textquotedblleft{}stone\textquotedblright{} (Middle Chinese \emph{dzyek}) and \emph{yu}, \textquotedblleft{}oil\textquotedblright{} (\emph{yuw}) (KANJIDIC2; Middle Chinese in Baxter's transcription).
\end{cousinweb}

\begin{grammarlens}[title={What you've built}]
One more word for giving reasons and opinions.
\end{grammarlens}

\subsection*{Your turn}
\begin{itemize}
\raggedright
  \item \emph{Say it:} \emph{sekiyu}
  \item \emph{Say it:} \emph{sekiyu}, once more
  \item \emph{Say it:} \emph{sekiyu wa sukunaku naru to omoimasu} — I think oil will become scarce.
  \item \emph{Recall:} say the Japanese for electric power, then the Japanese for fuel, then say \emph{sekiyu} again
\end{itemize}

\begin{tcolorbox}[breakable,colback=teal!4,colframe=teal!35!black,title={Before you move on}]
What does \ja{せきゆ} mean? (\textquotedblleft{}Petroleum; oil\textquotedblright{}.) Say it once more.
\end{tcolorbox}
\section[\texorpdfstring{\ja{もったいない} (mottainai)}{もったいない (mottainai)}]{\ja{もったいない} (mottainai) — wasteful; too good to waste}

\label{lesson:JA-C337-mottainai}



\begin{quote}
Before the new one: say the Japanese for fuel, then the Japanese for petroleum.
\end{quote}

\subsection*{You'll want to know: \ja{もったいない}}
\textbf{\ja{もったいない}} — \emph{mottainai} — \textquotedblleft{}wasteful; too good to waste\textquotedblright{}.

\begin{quote}
\textbf{\ja{もったいないから、すてないで} \ja{ください。}} — \emph{mottainai kara, sutenaide kudasai} — \textquotedblleft{}It's a waste, so please do not throw it away.\textquotedblright{}
\end{quote}

An i-adjective.

\begin{cousinweb}
\textbf{\ja{もったいない}} (\emph{mottainai}) is built from the Sino-Japanese \emph{mottai}, two Chinese characters read with their on'yomi, readings borrowed from Chinese, \emph{mochi}, \textquotedblleft{}not\textquotedblright{}, shortened to \emph{mot-}, and \emph{tai}, \textquotedblleft{}body\textquotedblright{} (Middle Chinese \emph{thejX}), and the native \emph{nai}, \textquotedblleft{}not\textquotedblright{}. JMdict notes that it is usually written in kana (KANJIDIC2; JMdict; Middle Chinese in Baxter's transcription).
\end{cousinweb}

\begin{grammarlens}[title={What you've built}]
One more word for giving reasons and opinions.
\end{grammarlens}

\subsection*{Your turn}
\begin{itemize}
\raggedright
  \item \emph{Say it:} \emph{mottainai}
  \item \emph{Say it:} \emph{mottainai}, once more
  \item \emph{Say it:} \emph{mottainai kara, sutenaide kudasai} — It's a waste, so please do not throw it away.
  \item \emph{Recall:} say the Japanese for fuel, then the Japanese for petroleum, then say \emph{mottainai} again
\end{itemize}

\begin{tcolorbox}[breakable,colback=teal!4,colframe=teal!35!black,title={Before you move on}]
What does \ja{もったいない} mean? (\textquotedblleft{}Wasteful; too good to waste\textquotedblright{}.) Say it once more.
\end{tcolorbox}
